\subsection{Why Innovation Succeeds Where Gradient Difference Fails}

The fundamental issue with gradient difference methods like Adan is \textbf{comparing two noisy samples}:
\begin{equation}
d_t = \g_t - \g_{t-1} = (\nabla f + \xi_t) - (\nabla f + \xi_{t-1}) = \underbrace{\text{signal}}_{\nabla f_t - \nabla f_{t-1}} + \underbrace{\text{noise}}_{\xi_t - \xi_{t-1}}
\end{equation}

The noise term has variance $\Var[\xi_t - \xi_{t-1}] \approx 2\sigma^2$, overwhelming the signal. No amount of variance smoothing ($\beta_3$) can fix this---the underlying correction signal is fundamentally noisy.

IRIS's innovation compares to \emph{smooth estimate}:
\begin{equation}
\Innovation_{t|t-1} = \g_t - \ghat_{t-1} = (\nabla f + \xi_t) - \nabla f_{\text{smooth}} = \text{signal} + \underbrace{\xi_t}_{\sigma^2}
\end{equation}

The noise is halved, enabling high variance smoothing ($\beta_3 = 0.999$) to produce hyperstable convergence.

\subsection{The Convergence-Divergence Problem}

Adan's most critical failure mode is \textbf{convergence followed by catastrophic divergence}:
\begin{itemize}
    \item Rosenbrock: Converges to $10^{-4}$ by step 200, then spikes to $10^{-1}$ at step 400 (10,000$\times$ increase)
    \item Rastrigin: Repeated spikes of $10^6\times$ at steps 600, 800
\end{itemize}

\textbf{Why this happens:} Near minima, $\g_t \approx \g_{t-1} \approx 0 \implies d_t \approx 0$. Correction vanishes while momentum remains large. Small perturbations cause unchecked momentum explosion.

\textbf{Why IRIS avoids this:} Innovation $\Innovation_{t|t-1} = \g_t - \ghat_{t-1}$ remains informative. Near minima, $\g_t \to 0$ but $\ghat_{t-1} > 0$ (lagging) creates \emph{large negative innovation} that actively damps updates.

This makes Adan unusable for:
\begin{itemize}
    \item \textbf{Checkpointing:} May save model during spike
    \item \textbf{Early stopping:} May stop during recovery phase
    \item \textbf{LR scheduling:} Spikes confuse schedulers
    \item \textbf{Distributed training:} Gradient explosions cause synchronization failures
\end{itemize}

IRIS's unconditional stability eliminates all these failure modes.

\subsection{Practical Implications}

\textbf{Hyperparameter tuning elimination:} IRIS's extreme LR robustness (works across LR $\in [0.3, 1.0]$) means practitioners can:
\begin{enumerate}
    \item Use LR = 1.0 as universal default
    \item Skip grid search entirely
    \item Transfer learning rates across tasks
\end{enumerate}

This reduces experimentation cost from $O(10)$ runs (grid search) to $O(1)$ (single run).

\textbf{No warmup needed:} Recursive bias correction provides exact correction from step 1, eliminating the 10-20 epoch warmup required by Adam/AdamW/Adan.

\textbf{Training acceleration:} CIFAR-100 speedup (113 vs 150 min) comes from:
\begin{itemize}
    \item Innovation reuse (compute once, use 3 times)
    \item Fused kernels (\texttt{lerp\_()})
    \item No GPU-CPU sync (recursive bias correction)
\end{itemize}

\subsection{Limitations}

\begin{itemize}
    \item \textbf{Additional hyperparameter:} $\beta_2$ must be tuned relative to $\beta_1$ (recommended: $\beta_2 = \beta_1 - 0.2$). However, this is \emph{far less sensitive} than learning rate.
    
    \item \textbf{Not always best:} On carefully-tuned Adam+schedule baselines, IRIS may match but not exceed performance. IRIS's advantage is \emph{eliminating the need for careful tuning}.
    
    \item \textbf{Theory-practice gap:} While we provide convergence guarantees, the extreme LR tolerance (1000$\times$ vs Adam) exceeds what theory currently predicts. Tighter bounds remain open.
\end{itemize}

\subsection{Future Work}

\begin{itemize}
    \item \textbf{Adaptive $\beta$ scheduling:} Can $\beta_1$ increase during training (Kalman-inspired)?
    \item \textbf{Layer-wise adaptation:} Different $\beta$ values per layer?
    \item \textbf{Second-order extensions:} Combine innovation with Hessian information?
    \item \textbf{Theoretical tightening:} Prove 1000$\times$ LR advantage formally
\end{itemize}
